\documentclass[UTF8,a4paper,10pt,openany,fontset=fandol]{ctexbook}
\usepackage[margin=18mm,headheight=14pt]{geometry}
\usepackage{amsmath,amssymb,booktabs,longtable,array}
\usepackage{xcolor,listings,fancyhdr,imakeidx,needspace}
\usepackage[unicode,colorlinks=true,linkcolor=blue!45!black,urlcolor=blue!45!black]{hyperref}
\makeindex[title=算法索引,columns=1]
\lstset{language=C++,basicstyle=\ttfamily\fontsize{8.5}{10.2}\selectfont,keywordstyle=\bfseries,
commentstyle=\color{black!60},numbers=left,numberstyle=\tiny\color{black!50},
numbersep=6pt,breaklines=true,breakatwhitespace=false,columns=fullflexible,
keepspaces=true,showstringspaces=false,tabsize=4,frame=single,rulecolor=\color{black!20}}
% Keep the last ten source lines together when a listing crosses a page.
\makeatletter
\newcount\CpcLastKeepLine
\lst@AddToHook{Init}{%
  \CpcLastKeepLine=\lst@lastline
  \advance\CpcLastKeepLine by -10
}
\lst@AddToHook{EveryPar}{%
  \ifnum\lst@lineno>\CpcLastKeepLine
    \vadjust{\nobreak}%
  \fi
}
\makeatother
\pagestyle{fancy}\fancyhf{}
\fancyhead[L]{\small Morning Flower and Evening Oath}
\fancyhead[R]{\small XCPC Template · Working Draft}
\fancyfoot[C]{\thepage}
\setcounter{tocdepth}{1}
\let\cleardoublepage\clearpage
\ctexset{chapter={beforeskip=8pt,afterskip=16pt}}
\begin{document}
\raggedbottom
\hypersetup{pageanchor=false}
\begin{titlepage}
\centering\vspace*{28mm}
{\Huge\bfseries XCPC Template Library\par}
\vspace{14mm}
{\LARGE Morning Flower\par and Evening Oath\par}
\vspace{12mm}{\large Tongji University\par}
\vspace{24mm}{\Large 双码风算法模板与数学速查\par}
\vspace{10mm}{\large 建设稿 · 2026-09-12\par}
\vfill
\begin{minipage}{0.82\textwidth}
\textbf{状态说明}\par
本稿包含首批已进行本地验证的模块与完整来源覆盖表的入口。
尚未完成 kuangbin、WIDA 全部条目的双码风迁移，尚未完成近三年中国赛站逐题覆盖审计，最大流、平衡树、双连通分量、可持久化等部分接口已有在线 AC 记录，具体题号、源码摘要及验证范围见 \texttt{verification/oj.json}。其余接口须分别核验，不能由同文件中某算法通过推断全部通过。
本地对拍通过不代表对所有输入的绝对正确性证明；适用范围、数值上限与退化约定属于接口的一部分。
\end{minipage}
\end{titlepage}
\hypersetup{pageanchor=true}
\frontmatter\tableofcontents
\chapter{使用与验证}
英文署名已于 QOJ 队伍账号 \texttt{ucup-team7502} 核对。
两套风格依据训练提交 2932575、2930805、2931279 提取，不擅自将共享队号的提交归属到个人。

代码采用 C++20。以赛时快速抄写为先，一行一条语句，函数体展开。无状态的简单算法使用独立函数，有状态算法使用轻量 struct。动态版默认需要 \texttt{bits/stdc++.h} 和 \texttt{cassert}；
传统版各头文件列明依赖。抄写单个函数或结构体时需同时抄写其依赖结构。
同名的不同风格实现通常用于不同提交，不建议混在同一编译单元中。
点号除特别注明外为 1-based，字符串和多项式为 0-based；所有区间约定应以条目说明为准。

测试源码位于 \texttt{tests/}，固定种子便于复现。
验证包含独立暴力参考、边界案例、静态与动态实现、AddressSanitizer 和 UndefinedBehaviorSanitizer。
\texttt{verification/} 保存实际运行记录与源码摘要；OJ 状态只在获得可查提交结果后登记。最大流记录为 \url{https://www.luogu.com.cn/record/297501480} 与 \url{https://www.luogu.com.cn/record/297502364}。

原始来源目录共登记 713 项，见 \texttt{docs/coverage.csv}。
它包含不同版本、例题及公式，不应将其数目视为独立算法数，也不能将“已登记”视为“已完成”。

\mainmatter
\part{紧凑动态风格}

\chapter{数据结构}

\Needspace{280pt}

\section{并查集}\label{compact-DSU}\index{DSU}

init(n)、find(x)、merge(x,y)，点号 1..n；siz[find(x)] 为集合大小。路径压缩与按大小合并，均摊 O($\alpha$(n))。

\lstinputlisting[firstline=6,lastline=28]{../src/compact/data_structure.hpp}

\Needspace{323pt}

\section{可撤销并查集}\label{compact-RollbackDSU}\index{RollbackDSU}

构造后 snapshot() 保存历史长度，rollback(t) 回退。禁止路径压缩；find/merge 为 O(log n)，回退每次成功合并 O(1)。

\lstinputlisting[firstline=29,lastline=56]{../src/compact/data_structure.hpp}

\Needspace{255pt}

\section{树状数组与第 k 小}\label{compact-Fenwick}\index{Fenwick}

add(x,v)、sum(x)、query(l,r)，下标 1..n，闭区间；kth(k) 要求所有频率非负且 k>0，超过总和返回 n+1。单次 O(log n)。

\lstinputlisting[firstline=57,lastline=76]{../src/compact/data_structure.hpp}

\Needspace{358pt}

\section{区间加与区间和}\label{compact-LazySeg}\index{LazySeg}

初值全零；add(l,r,v)、query(l,r)，1$\le$l$\le$r$\le$n。维护 sum 与加法 tag；所有中间和与长度乘积须在 long long 范围内。单次 O(log n)。

\lstinputlisting[firstline=77,lastline=108]{../src/compact/data_structure.hpp}

\Needspace{210pt}

\section{64 位异或线性基}\label{compact-XorBasis}\index{XorBasis}

insert(x) 返回是否增加秩；query(x) 返回 x 与所张成空间异或后的最大无符号值。单次 O(64)，不提供第 k 小。

\lstinputlisting[firstline=109,lastline=123]{../src/compact/data_structure.hpp}

\chapter{网络流}

\section{最大流与方案}\label{compact-Dinic}\index{Dinic}

Dinic(n)，add(u,v,c) 返回偶数原边编号；flow(s,t,limit) 返回本次新增流；used(id) 查询原边净流，cut(s) 返回残量可达点集。允许重边、自环，容量非负，s$\ne$t；重复求流须使用相同源汇。只有求尽最大流后 cut 才是最小割。一般 O(V$^{2}$E)，递归增广深度 O(V)。总流不得超过 long long 上限。

\lstinputlisting[firstline=6,lastline=67]{../src/compact/flow.hpp}

\section{势能最小费用最大流}\label{compact-MinCostFlow}\index{MinCostFlow}

add(u,v,cap,cost) 返回边编号，flow(s,t,limit) 返回新增流量及 128 位费用，used(id) 返回原边流量。允许负费用，初始残量网络必须没有负环；检测到负环抛异常。Bellman-Ford 初始化 O(VE)，每次增广 O(E log V)；增广次数不能简单假设为 O(V)。 另要求 $\sum$(cap×|cost|)<2\textasciicircum{}120，防止总费用超过 int128。

\lstinputlisting[firstline=68,lastline=131]{../src/compact/flow.hpp}

\Needspace{286pt}

\section{上下界可行循环流}\label{compact-BoundedCirculation}\index{BoundedCirculation}

add(u,v,lo,hi)，0$\le$lo$\le$hi；solve() 仅调用一次，返回可行性，仅成功后读取 used(id)。约束为每点流量守恒；不是带源汇最大流。所有上下界及其总和须在 long long 范围内。依赖 Dinic。

\lstinputlisting[firstline=132,lastline=154]{../src/compact/flow.hpp}

\chapter{图论}

\Needspace{344pt}

\section{非负权最短路}\label{compact-Dijkstra}\index{Dijkstra}

add(u,v,w)、run(s)，dis 为距离，pre 为前驱，path(t) 返回点序列。不可达用 LLONG\_MAX；所有可表示的最短距离必须严格小于该哨兵，w$\ge$0。O((V+E)log V)。

\lstinputlisting[firstline=6,lastline=35]{../src/compact/graph.hpp}

\Needspace{400pt}

\section{强连通分量}\label{compact-SCC}\index{SCC}

add(u,v)、run()，bel[1..n] 为分量编号，cnt 为数量；分量编号按缩点图拓扑序递增。迭代 Kosaraju 避免递归爆栈；O(V+E)，同时存正反图。

\lstinputlisting[firstline=36,lastline=72]{../src/compact/graph.hpp}

\Needspace{248pt}

\section{2-SAT 与任意可行方案}\label{compact-TwoSAT}\index{TwoSAT}

add(x,a,y,b) 添加 (x=a) 或 (y=b)，变量 1..n；solve() 返回可满足性，仅成功时读取 ans。依赖同文件 SCC；O(V+E)，不保证字典序最小。

\lstinputlisting[firstline=73,lastline=91]{../src/compact/graph.hpp}

\Needspace{381pt}

\section{二分图最大匹配}\label{compact-BipartiteMatching}\index{BipartiteMatching}

左右两侧分别编号 1..n、1..m；add(u,v)，solve() 返回匹配数，l[u] 与 r[v] 为匹配端点，0 表示未匹配。分层增广；本版本保守按 O(VE) 估计，不宣称已经验证 Hopcroft-Karp 最紧复杂度。

\lstinputlisting[firstline=92,lastline=125]{../src/compact/graph.hpp}

\Needspace{334pt}

\section{割点与桥}\label{compact-Lowlink}\index{Lowlink}

add(u,v) 返回无向边编号；run() 后 cut[u] 标识割点、bridge[id] 标识桥。按边编号跳过父边，支持重边、自环、不连通图。O(V+E)，递归深度可能 O(V)。

\lstinputlisting[firstline=126,lastline=154]{../src/compact/graph.hpp}

\chapter{树上算法}

\Needspace{555pt}

\section{重链剖分与 LCA}\label{compact-HLD}\index{HLD}

add(u,v) 建无向树；build(root) 后读取 fa、dep、siz、son、top、dfn、rk。子树对应 [dfn[u],dfn[u]+siz[u]-1]。lca(u,v) 为 O(log n)，path(u,v,work,edge) 分解路径；work(l,r) 仅适用于可交换聚合，边权模式排除 LCA 位置。输入必须是连通无环树，建树 O(n)，无递归。

\lstinputlisting[firstline=6,lastline=58]{../src/compact/tree.hpp}

\chapter{字符串}

\section{KMP、Z、Manacher 与最小表示}\label{compact-StringAlgo}\index{StringAlgo}

字符串及返回位置均 0-based。prefix(s)、z(s)；match(s,t) 查非空模式串并返回重叠匹配起点。manacher 返回奇数半径（含中心）和偶数半径（中心在 i 前）；rotation 返回最小循环表示起点，空串返回 0。各 O(n+m)。

\lstinputlisting[firstline=6,lastline=77]{../src/compact/string.hpp}

\Needspace{424pt}

\section{AC 自动机出现次数}\label{compact-AhoCorasick}\index{AhoCorasick}

字母表 a..z，add(s) 返回终止状态，要求模式非空；全部插入后 build()，只允许一次构建。count(text) 返回每个状态对应模式的出现次数，可重复查询；重复模式共享状态。构建 O(26N)，查询 O(|text|+N)。

\lstinputlisting[firstline=78,lastline=116]{../src/compact/string.hpp}

\Needspace{320pt}

\section{倍增后缀数组与 LCP}\label{compact-SuffixArray}\index{SuffixArray}

sa[k] 为第 k 个后缀起点，rk[i] 为排名，lcp[k] 为 sa[k-1] 与 sa[k] 的 LCP（lcp[0]=0）。空串合法；当前排序比较版 O(n log$^{2}$ n)，不能替代待补的线性 SA/DC3 条目。

\lstinputlisting[firstline=117,lastline=143]{../src/compact/string.hpp}

\Needspace{432pt}

\section{后缀自动机}\label{compact-SuffixAutomaton}\index{SuffixAutomaton}

extend(c) 接收整数 0..25；distinct() 返回不同非空子串数，counts() 返回 endpos 次数且不修改结构，可重复调用。状态 a，初始状态 0，link 根为 -1；空间 O(26n)，构建 O(26n)。

\lstinputlisting[firstline=144,lastline=183]{../src/compact/string.hpp}

\chapter{数论}

\section{模运算、素性、CRT 与类欧几里德}\label{compact-NumberTheory}\index{NumberTheory}

mul/power 支持无符号 64 位正模数，乘积用 uint128。prime 使用确定性 64 位底数。inverse(a,m) 无逆元返回 -1。crt(r,m,b,n) 合并同余式，无解返回 false，模数溢出抛异常。floor\_sum 为 i=0..n-1 的 floor((ai+b)/m) 之和，支持负 a、b，返回 int128；约定 n$\le$10$^{9}$、m$\le$10$^{9}$，a、b 为有符号 64 位。

\lstinputlisting[firstline=6,lastline=74]{../src/compact/number_theory.hpp}

\Needspace{327pt}

\section{Pollard–Rho 因数分解}\label{compact-PollardRho}\index{PollardRho}

factor(n) 返回升序质因子（含重数），n$\ge$1；依赖 NumberTheory。返回因子经过整除和素性检查；Las Vegas 随机运行时间，固定默认种子方便复现，可由调用者更换。无确定性最坏时间保证。

\lstinputlisting[firstline=75,lastline=102]{../src/compact/number_theory.hpp}

\Needspace{252pt}

\section{线性筛、欧拉函数与 Möbius 函数}\label{compact-LinearSieve}\index{LinearSieve}

构造参数 n$\ge$0，prime 为质数表，lp/phi/mu 为 0..n 数组。phi[1]=mu[1]=1，O(n) 时间和空间。

\lstinputlisting[firstline=103,lastline=122]{../src/compact/number_theory.hpp}

\Needspace{196pt}

\section{静态模整数}\label{compact-ModInt}\index{ModInt}

模板参数必须为不超过 INT\_MAX 的素数且$\ge$2。加减乘、pow、inv、除法；求逆要求非零，pow 指数非负。复杂度：加减乘 O(1)，幂和逆 O(log mod)。

\lstinputlisting[firstline=123,lastline=135]{../src/compact/number_theory.hpp}

\Needspace{199pt}

\section{阶乘组合数}\label{compact-Binomial}\index{Binomial}

Binomial<prime>(n)，要求 0$\le$n<prime；choose(n,k) 越界 k 返回 0。预处理 O(n+log prime)，单次 O(1)。尚不能替代合数模数 exLucas。

\lstinputlisting[firstline=136,lastline=148]{../src/compact/number_theory.hpp}

\chapter{多项式}

\section{NTT、形式幂级数、FWT 与递推}\label{compact-Polynomial}\index{Polynomial}

依赖 number\_theory.hpp。模数固定 998244353，乘法结果长度$\le$2$^{23}$；inverse 要求 a[0]$\ne$0，log 要求 a[0]=1，exp 要求 a[0]=0，截断长度 n>0。朴素逐层 NTT 实现的 FPS 为 O(n log n)。FWT 的 op 为 \textasciicircum{}、|、\&；BM 返回 s[n]=$\sum$c[i-1]s[n-i] 的系数；recurrence 用 O(k$^{2}$ log n) 求第 n 项。BM 对有限样本给出的递推不自动保证未知后续项成立。

\lstinputlisting[firstline=5,lastline=129]{../src/compact/polynomial.hpp}

\chapter{代数与数论进阶}

\section{模高斯消元、行列式、矩阵树定理}\label{compact-LinearAlgebra}\index{LinearAlgebra}

素数模数。solve(a,n) 输入增广矩阵，返回 consistent、rank、特解 particular 和零空间基 kernel。determinant(empty)=1。spanning\_trees(n,edges) 使用 0-based 无向图，允许重边，忽略自环；返回生成树数模 p。高斯与行列式 O(n$^{3}$)，矩阵幂 O(n$^{3}$ log e)。

\lstinputlisting[firstline=4,lastline=79]{../src/compact/algebra.hpp}

\Needspace{329pt}

\section{杜教筛}\label{compact-DuJiao}\index{DuJiao}

构造预筛上限 B$\ge$1，mertens(n) 返回 $\mu$ 的前缀和，totient\_sum(n) 返回 $\varphi$ 的前缀和（int128）。约定 0$\le$n$\le$10$^{11}$，按整除分块递归并缓存；大输入需要合理 B，不能把常数 B 当作性能保证。

\lstinputlisting[firstline=80,lastline=107]{../src/compact/algebra.hpp}

\Needspace{325pt}

\section{扩展 BSGS}\label{compact-DiscreteLog}\index{DiscreteLog}

solve(a,b,m) 返回最小 x$\ge$0 使 a\textasciicircum{}x$\equiv$b (mod m)，无解 -1；1$\le$m$\le$10$^{12}$，0\textasciicircum{}0=1。时间与内存 O($\sqrt{\vphantom x}$m)，unordered\_map 复杂度为期望界。

\lstinputlisting[firstline=108,lastline=135]{../src/compact/algebra.hpp}

\chapter{计算几何}

\section{整数精确二维几何}\label{compact-IntegerGeometry}\index{IntegerGeometry}

坐标绝对值$\le$10$^{12}$，点数$\le$10$^{6}$。叉积、点积、线段相交、凸包、点在多边形、点在凸包、旋转卡壳直径。凸包逆时针，不保留共线内点，不重复首点。contains 返回 0 外部、1 边界、2 内部。convex\_contains/diameter2 输入必须为本模板 strict hull。PolarLess 要求排除零向量。

\lstinputlisting[firstline=6,lastline=109]{../src/compact/geometry.hpp}

\section{浮点点线圆构造}\label{compact-RealGeometry}\index{RealGeometry}

long double，显式 eps 默认 10$^{-12}$。要求有限输入及适当尺度；接近退化时按容差分类，不能提供任意实数输入的精确拓扑保证。Result 区分 none/one/two/infinite/degenerate。零长度线返回 degenerate；投影与点到线段距离将零线段视为点。坐标缩放和误差预算应按具体题面复核。

\lstinputlisting[firstline=110,lastline=189]{../src/compact/geometry.hpp}

\chapter{优化与可持久化}

\Needspace{287pt}

\section{Max-plus 矩阵幂}\label{compact-MaxPlusMatrix}\index{MaxPlusMatrix}

a[i][j] 为转移收益，不存在的边为 neg；power(k) 求恰好 k 步的最优转移，单位矩阵对角为 0。所有有限中间值绝对值小于 2\textasciicircum{}119。O(n$^{3}$ log k)；2025 成都 B 的转移类型，具体题目的建模仍需调用者完成。

\lstinputlisting[firstline=6,lastline=28]{../src/compact/optimization.hpp}

\Needspace{492pt}

\section{离散李超树与最优直线编号}\label{compact-LiChao}\index{LiChao}

构造传入所有可能查询的横坐标，add(\{k,b,id\}) 插入直线，query(x) 返回最小值与最小编号，要求 x 在预登记坐标内且至少有一条线。整数求值用 int128，k、b、x 为有符号 64 位且中间值可表示。单次 O(log n)。

\lstinputlisting[firstline=29,lastline=75]{../src/compact/optimization.hpp}

\Needspace{386pt}

\section{主席树：静态区间第 k 小}\label{compact-PersistentKth}\index{PersistentKth}

构造传入 0-based 数组，kth(l,r,k) 的区间使用 1-based 闭区间；要求 1$\le$k$\le$r-l+1。离散化后每个前缀一个版本，单次 O(log n)，总空间 O(n log n)。

\lstinputlisting[firstline=76,lastline=109]{../src/compact/optimization.hpp}

\chapter{图论进阶}

\Needspace{389pt}

\section{矩形二分图最小权匹配}\label{compact-Hungarian}\index{Hungarian}

solve(cost) 输入 n×m 完整权值矩阵，0$\le$n$\le$m；每行匹配不同列，返回最小费用及列号（0-based）。允许负费用，使用 int128 势能；O(n$^{2}$m)。禁止边不能用未经证明的有限大常数随意替代，当前接口仅表示完整二分图。

\lstinputlisting[firstline=6,lastline=40]{../src/compact/graph_advanced.hpp}

\Needspace{390pt}

\section{有向最小树形图}\label{compact-Arborescence}\index{Arborescence}

solve(n,root,edges)，点号 0..n-1，n$\ge$1；返回以 root 为根向外可达的最小树形图费用，无解 nullopt。允许负权、重边、自环（自环不选）；O(VE)。当前仅返回费用，不恢复边方案。所有中间值须远小于 2\textasciicircum{}120。

\lstinputlisting[firstline=41,lastline=75]{../src/compact/graph_advanced.hpp}

\Needspace{340pt}

\section{无向全局最小割}\label{compact-StoerWagner}\index{StoerWagner}

solve(w)，n$\ge$2；w 是对称非负权矩阵，重边需先累加，主对角置零。返回割权与一侧顶点编号（0-based），允许不连通图。O(n$^{3}$)，总权值小于 2\textasciicircum{}120。

\lstinputlisting[firstline=76,lastline=105]{../src/compact/graph_advanced.hpp}

\chapter{精确几何进阶}

\section{整数最近点对与 Minkowski 和}\label{compact-GeometryExtra}\index{GeometryExtra}

closest\_pair(points) 返回最短距离平方，少于两个点为 nullopt，重复点返回 0，O(n log n)，全程精确整数。minkowski 输入 strict CCW 凸包，返回不重复首点的凸包；非退化 O(n+m)，退化时 O((n+m)log(n+m))。坐标及坐标和须满足整数几何范围。

\lstinputlisting[firstline=4,lastline=68]{../src/compact/geometry_extra.hpp}

\Needspace{233pt}

\section{整数三维共线、共面与方向}\label{compact-IntegerGeometry3D}\index{IntegerGeometry3D}

坐标绝对值$\le$10$^{9}$；orient(a,b,c,d) 为有向六倍四面体体积，0 表示共面；collinear 判三点共线，on\_segment 支持退化线段。int128 精确判定，不包含三维凸包构造。

\lstinputlisting[firstline=69,lastline=85]{../src/compact/geometry_extra.hpp}

\part{传统接口风格}

\chapter{数据结构}

\Needspace{365pt}

\section{并查集}\label{classic-DSU}\index{Disjoint\_Set}

算法约束见紧凑版本第~\pageref{compact-DSU}~页。本节代码独立可抄；采用下列实际接口名。

先声明容量模板参数，再调用 Init(n) 初始化；大对象必须置于全局或 static 存储。

\lstinputlisting[firstline=10,lastline=42]{../src/classic/data_structure.hpp}

\Needspace{382pt}

\section{可撤销并查集}\label{classic-RollbackDSU}\index{Rollback\_DSU}

算法约束见紧凑版本第~\pageref{compact-RollbackDSU}~页。本节代码独立可抄；采用下列实际接口名。

\lstinputlisting[firstline=123,lastline=157]{../src/classic/data_structure.hpp}

\Needspace{374pt}

\section{树状数组与第 k 小}\label{classic-Fenwick}\index{Binary\_Indexed\_Tree}

算法约束见紧凑版本第~\pageref{compact-Fenwick}~页。本节代码独立可抄；采用下列实际接口名。

先声明容量模板参数，再调用 Init(n) 初始化；大对象必须置于全局或 static 存储。

\lstinputlisting[firstline=43,lastline=76]{../src/classic/data_structure.hpp}

\Needspace{477pt}

\section{区间加与区间和}\label{classic-LazySeg}\index{Segment\_Tree}

算法约束见紧凑版本第~\pageref{compact-LazySeg}~页。本节代码独立可抄；采用下列实际接口名。

先声明容量模板参数，再调用 Init(n) 初始化；大对象必须置于全局或 static 存储。

\lstinputlisting[firstline=77,lastline=122]{../src/classic/data_structure.hpp}

\Needspace{244pt}

\section{64 位异或线性基}\label{classic-XorBasis}\index{Xor\_Basis}

算法约束见紧凑版本第~\pageref{compact-XorBasis}~页。本节代码独立可抄；采用下列实际接口名。

\lstinputlisting[firstline=158,lastline=176]{../src/classic/data_structure.hpp}

\chapter{网络流}

\section{最大流与方案}\label{classic-Dinic}\index{Network\_Flow}

算法约束见紧凑版本第~\pageref{compact-Dinic}~页。本节代码独立可抄；采用下列实际接口名。

先声明容量模板参数，再调用 Init(n) 初始化；大对象必须置于全局或 static 存储。

\lstinputlisting[firstline=9,lastline=79]{../src/classic/flow.hpp}

\section{势能最小费用最大流}\label{classic-MinCostFlow}\index{Min\_Cost\_Flow}

算法约束见紧凑版本第~\pageref{compact-MinCostFlow}~页。本节代码独立可抄；采用下列实际接口名。

使用 Insert 加边、Flow 求流、Used 查看方案。

\lstinputlisting[firstline=21,lastline=117]{../src/classic/min_cost_flow.hpp}

\Needspace{414pt}

\section{上下界可行循环流}\label{classic-BoundedCirculation}\index{Bounded\_Circulation}

算法约束见紧凑版本第~\pageref{compact-BoundedCirculation}~页。本节代码独立可抄；采用下列实际接口名。

\lstinputlisting[firstline=80,lastline=117]{../src/classic/flow.hpp}

\chapter{图论}

\Needspace{497pt}

\section{非负权最短路}\label{classic-Dijkstra}\index{Shortest\_Path}

算法约束见紧凑版本第~\pageref{compact-Dijkstra}~页。本节代码独立可抄；采用下列实际接口名。

\lstinputlisting[firstline=22,lastline=69]{../src/classic/graph.hpp}

\section{强连通分量}\label{classic-SCC}\index{Strong\_Component}

算法约束见紧凑版本第~\pageref{compact-SCC}~页。本节代码独立可抄；采用下列实际接口名。

\lstinputlisting[firstline=70,lastline=132]{../src/classic/graph.hpp}

\Needspace{333pt}

\section{2-SAT 与任意可行方案}\label{classic-TwoSAT}\index{Two\_SAT}

算法约束见紧凑版本第~\pageref{compact-TwoSAT}~页。本节代码独立可抄；采用下列实际接口名。

\lstinputlisting[firstline=133,lastline=161]{../src/classic/graph.hpp}

\section{二分图最大匹配}\label{classic-BipartiteMatching}\index{Bipartite\_Matching}

算法约束见紧凑版本第~\pageref{compact-BipartiteMatching}~页。本节代码独立可抄；采用下列实际接口名。

\lstinputlisting[firstline=162,lastline=225]{../src/classic/graph.hpp}

\Needspace{504pt}

\section{割点与桥}\label{classic-Lowlink}\index{Low\_Link}

算法约束见紧凑版本第~\pageref{compact-Lowlink}~页。本节代码独立可抄；采用下列实际接口名。

\lstinputlisting[firstline=226,lastline=274]{../src/classic/graph.hpp}

\chapter{树上算法}

\section{重链剖分与 LCA}\label{classic-HLD}\index{Heavy\_Light\_Decomposition}

算法约束见紧凑版本第~\pageref{compact-HLD}~页。本节代码独立可抄；采用下列实际接口名。

\lstinputlisting[firstline=22,lastline=111]{../src/classic/tree.hpp}

\chapter{字符串}

\section{KMP、Z、Manacher 与最小表示}\label{classic-StringAlgo}\index{String\_Algorithm}

算法约束见紧凑版本第~\pageref{compact-StringAlgo}~页。本节代码独立可抄；采用下列实际接口名。

\lstinputlisting[firstline=22,lastline=136]{../src/classic/string.hpp}

\section{AC 自动机出现次数}\label{classic-AhoCorasick}\index{AC\_Automaton}

算法约束见紧凑版本第~\pageref{compact-AhoCorasick}~页。本节代码独立可抄；采用下列实际接口名。

\lstinputlisting[firstline=137,lastline=206]{../src/classic/string.hpp}

\Needspace{473pt}

\section{倍增后缀数组与 LCP}\label{classic-SuffixArray}\index{Suffix\_Array}

算法约束见紧凑版本第~\pageref{compact-SuffixArray}~页。本节代码独立可抄；采用下列实际接口名。

\lstinputlisting[firstline=207,lastline=251]{../src/classic/string.hpp}

\section{后缀自动机}\label{classic-SuffixAutomaton}\index{Suffix\_Automaton}

算法约束见紧凑版本第~\pageref{compact-SuffixAutomaton}~页。本节代码独立可抄；采用下列实际接口名。

\lstinputlisting[firstline=252,lastline=322]{../src/classic/string.hpp}

\chapter{数论}

\section{模运算、素性、CRT 与类欧几里德}\label{classic-NumberTheory}\index{Number\_Theory}

算法约束见紧凑版本第~\pageref{compact-NumberTheory}~页。本节代码独立可抄；采用下列实际接口名。

\lstinputlisting[firstline=22,lastline=155]{../src/classic/number_theory.hpp}

\Needspace{497pt}

\section{Pollard–Rho 因数分解}\label{classic-PollardRho}\index{Pollard\_Rho}

算法约束见紧凑版本第~\pageref{compact-PollardRho}~页。本节代码独立可抄；采用下列实际接口名。

\lstinputlisting[firstline=156,lastline=203]{../src/classic/number_theory.hpp}

\Needspace{380pt}

\section{线性筛、欧拉函数与 Möbius 函数}\label{classic-LinearSieve}\index{Linear\_Sieve}

算法约束见紧凑版本第~\pageref{compact-LinearSieve}~页。本节代码独立可抄；采用下列实际接口名。

\lstinputlisting[firstline=204,lastline=238]{../src/classic/number_theory.hpp}

\Needspace{392pt}

\section{静态模整数}\label{classic-ModInt}\index{Mod\_Int}

算法约束见紧凑版本第~\pageref{compact-ModInt}~页。本节代码独立可抄；采用下列实际接口名。

\lstinputlisting[firstline=239,lastline=274]{../src/classic/number_theory.hpp}

\Needspace{275pt}

\section{阶乘组合数}\label{classic-Binomial}\index{Combination}

算法约束见紧凑版本第~\pageref{compact-Binomial}~页。本节代码独立可抄；采用下列实际接口名。

\lstinputlisting[firstline=275,lastline=296]{../src/classic/number_theory.hpp}

\chapter{多项式}

\section{NTT、形式幂级数、FWT 与递推}\label{classic-Polynomial}\index{Polynomial}

算法约束见紧凑版本第~\pageref{compact-Polynomial}~页。本节代码独立可抄；采用下列实际接口名。

\lstinputlisting[firstline=5,lastline=224]{../src/classic/polynomial.hpp}

\chapter{代数与数论进阶}

\section{模高斯消元、行列式、矩阵树定理}\label{classic-LinearAlgebra}\index{Linear\_Algebra}

算法约束见紧凑版本第~\pageref{compact-LinearAlgebra}~页。本节代码独立可抄；采用下列实际接口名。

\lstinputlisting[firstline=4,lastline=132]{../src/classic/algebra.hpp}

\Needspace{499pt}

\section{杜教筛}\label{classic-DuJiao}\index{Du\_Jiao}

算法约束见紧凑版本第~\pageref{compact-DuJiao}~页。本节代码独立可抄；采用下列实际接口名。

\lstinputlisting[firstline=133,lastline=180]{../src/classic/algebra.hpp}

\Needspace{470pt}

\section{扩展 BSGS}\label{classic-DiscreteLog}\index{Discrete\_Log}

算法约束见紧凑版本第~\pageref{compact-DiscreteLog}~页。本节代码独立可抄；采用下列实际接口名。

\lstinputlisting[firstline=181,lastline=225]{../src/classic/algebra.hpp}

\chapter{计算几何}

\section{整数精确二维几何}\label{classic-IntegerGeometry}\index{Integer\_Geometry}

算法约束见紧凑版本第~\pageref{compact-IntegerGeometry}~页。本节代码独立可抄；采用下列实际接口名。

\lstinputlisting[firstline=22,lastline=183]{../src/classic/geometry.hpp}

\section{浮点点线圆构造}\label{classic-RealGeometry}\index{Real\_Geometry}

算法约束见紧凑版本第~\pageref{compact-RealGeometry}~页。本节代码独立可抄；采用下列实际接口名。

\lstinputlisting[firstline=184,lastline=332]{../src/classic/geometry.hpp}

\chapter{优化与可持久化}

\Needspace{381pt}

\section{Max-plus 矩阵幂}\label{classic-MaxPlusMatrix}\index{Max\_Plus\_Matrix}

算法约束见紧凑版本第~\pageref{compact-MaxPlusMatrix}~页。本节代码独立可抄；采用下列实际接口名。

\lstinputlisting[firstline=22,lastline=55]{../src/classic/optimization.hpp}

\section{离散李超树与最优直线编号}\label{classic-LiChao}\index{Li\_Chao\_Tree}

算法约束见紧凑版本第~\pageref{compact-LiChao}~页。本节代码独立可抄；采用下列实际接口名。

\lstinputlisting[firstline=56,lastline=146]{../src/classic/optimization.hpp}

\section{主席树：静态区间第 k 小}\label{classic-PersistentKth}\index{Persistent\_Kth}

算法约束见紧凑版本第~\pageref{compact-PersistentKth}~页。本节代码独立可抄；采用下列实际接口名。

\lstinputlisting[firstline=147,lastline=211]{../src/classic/optimization.hpp}

\chapter{图论进阶}

\section{矩形二分图最小权匹配}\label{classic-Hungarian}\index{Kuhn\_Munkres}

算法约束见紧凑版本第~\pageref{compact-Hungarian}~页。本节代码独立可抄；采用下列实际接口名。

\lstinputlisting[firstline=22,lastline=87]{../src/classic/graph_advanced.hpp}

\section{有向最小树形图}\label{classic-Arborescence}\index{Directed\_MST}

算法约束见紧凑版本第~\pageref{compact-Arborescence}~页。本节代码独立可抄；采用下列实际接口名。

\lstinputlisting[firstline=88,lastline=153]{../src/classic/graph_advanced.hpp}

\Needspace{519pt}

\section{无向全局最小割}\label{classic-StoerWagner}\index{Global\_Min\_Cut}

算法约束见紧凑版本第~\pageref{compact-StoerWagner}~页。本节代码独立可抄；采用下列实际接口名。

\lstinputlisting[firstline=154,lastline=204]{../src/classic/graph_advanced.hpp}

\chapter{精确几何进阶}

\section{整数最近点对与 Minkowski 和}\label{classic-GeometryExtra}\index{Geometry\_Extra}

算法约束见紧凑版本第~\pageref{compact-GeometryExtra}~页。本节代码独立可抄；采用下列实际接口名。

\lstinputlisting[firstline=4,lastline=103]{../src/classic/geometry_extra.hpp}

\Needspace{420pt}

\section{整数三维共线、共面与方向}\label{classic-IntegerGeometry3D}\index{Integer\_Geometry\_3D}

算法约束见紧凑版本第~\pageref{compact-IntegerGeometry3D}~页。本节代码独立可抄；采用下列实际接口名。

\lstinputlisting[firstline=104,lastline=142]{../src/classic/geometry_extra.hpp}

\part{数学速查与维护}
\chapter{数学速查：条件先于公式}
\section{整除、卷积与筛法}\index{Dirichlet convolution}\index{Mobius inversion}
对正整数定义狄利克雷卷积
\[
(f*g)(n)=\sum_{d\mid n}f(d)g(n/d).
\]
单位元为 $\varepsilon(n)=[n=1]$，常函数 $1(n)=1$，恒等函数 $\operatorname{id}(n)=n$。
关键恒等式是 $\mu*1=\varepsilon$、$\varphi*1=\operatorname{id}$。
若 $F(n)=\sum_{d\mid n}f(d)$，则 $f(n)=\sum_{d\mid n}\mu(d)F(n/d)$。
\[
\sum_{i=1}^{a}\sum_{j=1}^{b}[\gcd(i,j)=1]
=\sum_{d=1}^{\min(a,b)}\mu(d)\lfloor a/d\rfloor\lfloor b/d\rfloor.
\]
整除分块：当左端点为 $l$、商为 $q=\lfloor n/l\rfloor$，右端点为 $r=\lfloor n/q\rfloor$。
同时包含多个商时，取各自右端点的最小值。
杜教筛由卷积恒等式获得
\[
M(n)=1-\sum_{l=2}^{n}M(\lfloor n/l\rfloor),\qquad
\Phi(n)=\frac{n(n+1)}2-\sum_{l=2}^{n}\Phi(\lfloor n/l\rfloor),
\]
求和按等商区间合并；第二式求和中没有额外乘以 $l$。

\section{模运算的适用条件}\index{CRT}\index{Euler theorem}\index{Lucas theorem}
模逆元存在当且仅当 $\gcd(a,m)=1$。费马逆元 $a^{p-2}$ 仅适用于素数模数及 $a\not\equiv0$。
对正模数 $m,n$，同余式 $x\equiv a\pmod m$、$x\equiv b\pmod n$ 可合并当且仅当
$\gcd(m,n)\mid(b-a)$；新模数为 $\operatorname{lcm}(m,n)$。
\[
\binom nk\equiv\prod_i\binom{n_i}{k_i}\pmod p\quad(p\text{ 为素数，}n_i,k_i\text{ 是 }p\text{ 进制位}).
\]
Lucas 不能直接套在合数模数。阶乘与逆阶乘预处理要求上限小于模数；超过模数会出现阶乘为零。
扩展欧拉降幂中，非互质情况不能直接将指数对 $\varphi(m)$ 取模；需要判断原指数是否达到足够阈值。
大整数指数可以在读入时同时维护余数与是否超过阈值。

\section{素数模平方根}\index{Cipolla}\index{quadratic residue}
奇素数 $p$ 下，对 $n\not\equiv0\pmod p$，Euler 判别式
$n^{(p-1)/2}\equiv1\pmod p$ 当且仅当存在平方根，否则为 $-1$。
当 $p\equiv3\pmod4$ 且有解时，$r=n^{(p+1)/4}$ 为一根，另一根是 $p-r$。

Cipolla 选择 $a$ 使 $w=a^2-n$ 是非平方，令 $\omega^2=w$。
在扩域中取 $\alpha=a+\omega$，有 $\alpha^p=a-\omega$，故 $\alpha^{p+1}=n$。
令 $\beta=\alpha^{(p+1)/2}$，则 $\beta^2=n$ 且
$\beta^{p-1}=n^{(p-1)/2}=1$，所以 $\beta$ 落在原素数域中。
实现见第~\pageref{compact-mod_sqrt}~页；零和模数 2 单独处理。
模数为合数时不能照搬此判别和扩域构造。

\section{合数模数组合数与插值}\index{exLucas}\index{Lagrange interpolation}
对素数幂 $q=p^e$，记 $v_p(n!)=\sum_{j\ge1}\lfloor n/p^j\rfloor$，
$U_p(n)=n!/p^{v_p(n!)}$。设
$t=v_p(n!)-v_p(k!)-v_p((n-k)!)$。当 $t\ge e$ 时组合数模 $q$ 为零；否则
\[
\binom nk\equiv U_p(n)\,U_p(k)^{-1}\,U_p(n-k)^{-1}\,p^t\pmod q.
\]
这里仅对与 $p$ 互素的单位部分求逆。各素数幂上的结果通过 CRT 合并，
对应第~\pageref{compact-ExLucas}~页的实现。

在素数域中，$n$ 个模 $p$ 两两不同的横坐标确定唯一的次数小于 $n$ 的多项式。
定义 $w_i=y_i/\prod_{j\ne i}(x_i-x_j)$，则
\[
f(k)=\sum_{i=0}^{n-1}w_i\prod_{j<i}(k-x_j)\prod_{j>i}(k-x_j).
\]
该式不除以 $k-x_i$，查询点命中样本点时仍然成立。
当 $x_i=i$ 时，分母为 $i!\,(n-1-i)!\,(-1)^{n-1-i}$，要求 $n\le p$。
实现见第~\pageref{compact-Lagrange}~页。
对整数幂前缀和 $S_d(k)=\sum_{i=1}^k i^d$，当 $0\le d\le p-2$ 时，
它在模 $p$ 意义下可用次数至多 $d+1$ 的多项式表示；取 $0,1,\ldots,d+1$
的样本即可。不能只因某序列前几项能插值，就断言后续仍由该多项式给出。

\section{组合计数}\index{inclusion-exclusion}\index{Catalan}\index{Stirling numbers}
容斥：$|\bigcup A_i|=\sum_{\emptyset\ne S}(-1)^{|S|+1}|\bigcap_{i\in S}A_i|$。
二项式反演：若 $b_n=\sum_{k=0}^{n}\binom nk a_k$，则
$a_n=\sum_{k=0}^{n}(-1)^{n-k}\binom nk b_k$。
\[
\sum_k\binom ak\binom b{n-k}=\binom{a+b}n,\qquad
C_n=\binom{2n}n-\binom{2n}{n+1}.
\]
差分形式的 Catalan 公式避免直接除以 $n+1$，在模数下尤其有用，但组合数本身仍须正确求模。
第二类 Stirling 数满足
\[
\left\{\begin{matrix}n\\k\end{matrix}\right\}
=\left\{\begin{matrix}n-1\\k-1\end{matrix}\right\}
+k\left\{\begin{matrix}n-1\\k\end{matrix}\right\},
\quad S(0,0)=1.
\]
第一类无符号 Stirling 数把系数 $k$ 换为 $n-1$。
\begin{center}
\begin{tabular}{lll}
\toprule
球与盒 & 空盒允许 & 空盒禁止\\\midrule
不同球，不同盒 & $k^n$ & $k!S(n,k)$\\
相同球，不同盒 & $\binom{n+k-1}{k-1}$ & $\binom{n-1}{k-1}$\\
不同球，相同盒 & $\sum_{j=0}^{k}S(n,j)$ & $S(n,k)$\\
\bottomrule
\end{tabular}
\end{center}
上表默认 $k\ge1$；$n=0$、$k=0$ 等边界应按原问题约定单独处理。

\section{群作用计数}\index{Burnside}\index{Polya}
有限群 $G$ 作用在有限集合 $X$ 上时，轨道数为
\[
|X/G|=\frac1{|G|}\sum_{g\in G}|\operatorname{Fix}(g)|.
\]
对位置置换的 $q$ 色染色，$|\operatorname{Fix}(g)|=q^{c(g)}$，其中 $c(g)$ 为循环数。
环形项链只商去旋转；手链还商去反射，二者群不同，不能共用分母。
模数整除 $|G|$ 时不能使用模逆元，须先在合适的整数模数下求和再做精确除法。

\section{生成函数与递推}\index{generating functions}\index{Berlekamp-Massey}
普通生成函数 $A(x)=\sum a_nx^n$ 的乘法对应普通卷积。
指数生成函数 $A(x)=\sum a_nx^n/n!$ 的乘法对应
$c_n=\sum_k\binom nk a_kb_{n-k}$。
形式幂级数求逆要求常数项可逆；$\log A$ 取 $A(0)=1$，$\exp A$ 取 $A(0)=0$。
有限域中积分涉及 $1/i$，必须避免 $i$ 被特征整除。
\[
B_{\rm new}=B(2-AB)\pmod{x^{2m}},\qquad
(\log A)'=A'/A.
\]
BM 从样本得到最短线性递推。要对未来项使用，必须另外证明序列确实具有不超过相应阶数的线性递推；
少量数值吻合不是证明。矩阵转移、有限状态计数常可提供阶数上界。

\section{线性代数与生成树}\index{rank}\index{Matrix Tree theorem}
域上的线性方程组一致时，解集为一个特解加零空间，零空间维数为未知数个数减秩。
浮点消元中的秩依赖误差阈值；模素数消元是精确域运算；合数模数下不可随意除以非零元。
无向图 Laplacian 的对角元为度数，非对角元为负的边重数。删去同一行列后的行列式为生成树数。
自环不参与生成树；重边需要累加。只有一个顶点时空余子式行列式约定为 1。
有向图需要重新明确入度/出度 Laplacian 与根向内/向外树的方向，不能直接复用无向封装。

\section{概率、期望与博弈}\index{expectation}\index{Sprague-Grundy}
期望线性性不要求独立：$\mathbb E[\sum X_i]=\sum\mathbb E[X_i]$。
方差可加需要协方差为零；独立是充分条件。非负整数变量
$\mathbb E[X]=\sum_{k\ge1}\Pr(X\ge k)$。
吸收过程的期望方程 $E_i=1+\sum_jp_{ij}E_j$ 需要检查是否几乎必然吸收及期望是否有限；
不能对闭合非吸收分量直接求唯一有限解。
有限无环公平组合游戏在正常玩法下
\[
\operatorname{SG}(u)=\operatorname{mex}\{\operatorname{SG}(v):u\to v\}.
\]
独立子游戏的 SG 值异或为零当且仅当先手必败。
反常玩法、含环游戏、双方可走步不同的游戏不能直接套这个结论。

\section{数值算法与几何}\index{numerical stability}\index{Pick theorem}
排序必须使用严格弱序；不能以“差小于 eps 则相等”定义通用排序比较器。
叉积与行列式容易发生抵消；整数题优先精确整数判定，浮点构造单独处理误差。
对于简单格点多边形（无孔），Pick 定理 $A=I+B/2-1$，边界格点数
$B=\sum\gcd(|\Delta x|,|\Delta y|)$；不要重复计入边端点。
自适应积分的局部误差估计需要函数足够规则；奇点、间断与强振荡必须先分段处理。
三分搜索要求单峰性，整数域最后枚举剩余区间；二分答案要求判定单调性。

\paragraph{来源索引}
本章为重新组织的公式速查，扩展范围依据 \url{https://oi-wiki.org/math/}。
小模数 exLucas 与有限域拉格朗日插值已有双风格实现和本地验证；素数模平方根已有 Cipolla 实现；仍需补齐二次剩余的其他方法与定理、原根、Min25、Newton 插值、子集卷积、行列式进阶、数值积分等内容及验证。本章公式不能代替这些待完成代码。

\section{线性递推：系数方向与外推条件}
\index{线性递推}
约定 $a_i=\sum_{j=1}^{k}c_{j-1}a_{i-j}$，初值是
$a_0,\ldots,a_{k-1}$。对应特征多项式为
\[
F(x)=x^k-c_0x^{k-1}-c_1x^{k-2}-\cdots-c_{k-1}.
\]
若 $x^n\bmod F(x)=\sum_{i=0}^{k-1}r_ix^i$，则
$a_n=\sum_{i=0}^{k-1}r_ia_i$。代码由高次向低次消去多项式项，
第 $j$ 个系数写入原次数减 $j$ 的位置，避免倒置系数。
例如 Fibonacci 初值为 $[0,1]$，递推系数为 $[1,1]$。

BM 在域上求有限前缀的最短递推。本书实现要求素数模数；
一般的模合数环不能保证非零偏差存在逆元。
仅凭若干样本无法证明真实序列以后一直遵循同一个规律。
若事先知道序列满足阶数至多 $d$ 的齐次线性递推，
至少给出前 $2d$ 项，再用恢复出的系数求后续项。
系数向量末尾的零不能擅自删除：它可能对应非平稳的初始段，
改变阶数会改变递推开始适用的位置。

\chapter{计算几何：精确判定与适用范围}
\section{整数系数的有界半平面交}
实现见第~\pageref{compact-IntegerHalfplanes}~页。
输入是不等式 $ax+by\le c$，三个系数均为绝对值不超过 $10^9$ 的整数。
本接口要求交集有界，只输出正面积多边形；无解、点交集和线段交集都返回空向量。
它不是一般的可行性分类接口，不可将其空返回直接解读为闭半平面交集必为空。
无界输入不在约定范围内，也不偷偷添加一个人工边框。

两条非平行边界交点用齐次坐标 $(X,Y,D)$ 表示：
\[
D=a_1b_2-b_1a_2,\quad X=c_1b_2-b_1c_2,\quad Y=a_1c_2-c_1a_2.
\]
若 $D<0$，三者同时取负，使 $D>0$。
测试 $aX+bY\le cD$ 即可判断交点位于半平面内，无需除法。
系数界保证 $|X|,|Y|,D\le2\times10^{18}$，可用 64 位保存；
相等判断的交叉乘积至多为 $4\times10^{36}$，使用 128 位。
因此本接口内的构造、排序和位置判定都是精确整数运算。

按法向量的半圆与叉积排序，比较器不使用 epsilon。
同方向平行约束按正规化后的常数项保留更紧者，避免除以法向量长度。
扫描时若新约束排除队首或队尾交点，就删除相应无效边界；
每条边界最多入队、出队一次，排序后为线性时间。
最后闭合队列，清除重复顶点；有界的非正面积结果不会输出多边形。
零法向量表示常真约束或直接矛盾，先单独处理。

测试枚举所有边界对的交点，用任意精度整数逐条验证可行性，
将可行顶点集合与返回集合比较，并检查循环顺序的严格凸性。
两套版本通过 5000 组有界随机输入、平行和退化案例、近乎平行的十亿系数，
以及十万条冗余约束的规模测试；普通与 ASan/UBSan 均通过。
将有理数坐标转换为浮点、计算面积或处理浮点输入属于另一个数值问题，尚未由此验证。
WIDA 的一般半平面交条目保留部分完成状态，浮点版本与无界分类仍需补齐。

\clearpage
\section{最近点对、Minkowski 和与三维整数判定}
最近点对与 Minkowski 和见第~\pageref{compact-GeometryExtra}~页。
最近点对返回距离平方，不开方；少于两个点时返回空 optional，重复点返回零。
分治使用递归，每层按纵坐标归并，总时间 $O(n\log n)$，额外空间 $O(n)$。
坐标绝对值上限为 $10^{12}$，作差后平方使用 128 位。

Minkowski 和输入必须是严格逆时针凸包，不重复首点，不含边上的冗余共线点；
允许循环移动起点，也允许空集、单点和线段。
输入坐标以及任意两个输入点的坐标和都须在 $[-10^{12},10^{12}]$ 内。
非退化凸包按边方向归并，时间 $O(n+m)$；退化分支枚举点和后求凸包，
时间 $O((n+m)\log(n+m))$。
返回同样不重复首点的严格凸包。

两套代码通过任意精度整数距离枚举，以及对所有点和独立执行礼品包装法的对拍；
覆盖退化、循环起点和大坐标。规模测试包含二十万点最近点对，
以及 $40001$ 顶点的抛物线凸包与自身求和，结果应为原凸包各坐标的两倍。
这些属于本地验证，不表示新增在线 AC。

三维判定见第~\pageref{compact-IntegerGeometry3D}~页。
点的三个坐标为 \texttt{x,y,z}，绝对值均不超过 $10^9$。
\texttt{orient(a,b,c,d)} 返回有向六倍四面体体积，零表示四点共面，
交换任意两个点会改变符号。\texttt{collinear} 判断三点共线，
\texttt{on\_segment} 包含端点，且支持零长度线段。
传统版线段接口名为 \texttt{On\_Segment}。

点差每个分量至多 $2\times10^9$，叉积分量至多 $8\times10^{18}$，
三重积绝对值至多 $4.8\times10^{28}$，在有符号 128 位范围内。
公开的叉积、点积辅助函数不对任意 128 位向量保证不溢出；上述保证针对点差及其叉积。
任意精度整数按 $4\times4$ 行列式的 24 个排列独立计算方向，
并用二阶子式及坐标区间验证共线与在线段上。
已通过立方体八个顶点的全部有序四元组、随机十亿级坐标、换序及退化测试，
普通与 ASan/UBSan 均通过。这里不包含三维凸包构造。

\chapter{来源与覆盖维护}
\section{参考基线}
\begin{itemize}
\item kuangbin，ACM 常用算法模板，2018 年版本：\url{https://kuangbin.github.io/2018/08/01/ACM-template/}。
\item WIDA XCPC：\url{https://github.com/hh2048/XCPC}，参考快照 \texttt{484e9e0}。
\item OI Wiki 数学：\url{https://oi-wiki.org/math/}。
\end{itemize}
本项目重新实现的代码与公式速查不能冒用原作者的验证记录。
完整覆盖审计应为每个来源条目记录实现文件、替代关系与验证证据，保留遗漏项直至完成。
\section{后续覆盖项}
尚需完成的范围以覆盖表为准，包括动态树、平衡树、一般图匹配、最小树形图、上下界流、
高阶几何构造、半平面交、三维凸包、数值积分、复杂数论及多项式、搜索与动态规划等。
不能用通用模块的存在推断其覆盖了所有原书例题。
\backmatter\printindex
\end{document}
